\documentclass[11pt]{article}
\usepackage[utf8]{inputenc}
\usepackage{amsmath,amssymb,amsthm}
\usepackage{algorithm}
\usepackage{algpseudocode}
\usepackage{booktabs}
\usepackage{graphicx}
\usepackage{hyperref}
\usepackage{xcolor}
\usepackage[normalem]{ulem}  % For \sout strikethrough
\usepackage[margin=1in]{geometry}

\newtheorem{theorem}{Theorem}
\newtheorem{lemma}[theorem]{Lemma}
\newtheorem{proposition}[theorem]{Proposition}
\theoremstyle{definition}
\newtheorem{definition}{Definition}
\theoremstyle{remark}
\newtheorem{remark}{Remark}

\title{Global vs.\ Per-Comparison Dispersion and Size Factor Estimation\\in Single-Cell CRISPR Screens}
\author{crispyx Development Team}
\date{January 2026}

\begin{document}

\maketitle

\begin{abstract}
We present a computational innovation for negative binomial generalized linear model (NB-GLM) analysis of single-cell CRISPR perturbation screens. Traditional differential expression pipelines (DESeq2, PyDESeq2) estimate dispersion parameters \emph{per comparison}, which scales as $O(K)$ for $K$ perturbations. We introduce \textbf{global dispersion estimation} that computes the dispersion trend once using all cells, reducing complexity to $O(1)$ per perturbation. This document details the biological motivation, mathematical formulation, and empirical validation showing perfect correlation ($\rho = 1.000$) between global and per-comparison approaches with \textbf{1.7--2.3$\times$ speedup} on real datasets.
\end{abstract}

\tableofcontents
\newpage

%==============================================================================
\section{Introduction: The CRISPR Screen Analysis Challenge}
%==============================================================================

Single-cell CRISPR screens generate datasets with:
\begin{itemize}
    \item $N \sim 10^4$--$10^6$ cells
    \item $G \sim 5{,}000$--$20{,}000$ genes
    \item $K \sim 10^2$--$10^4$ perturbations (guide RNA targets)
\end{itemize}

For each perturbation $k$, we fit an NB-GLM comparing treated cells to controls:
\[
Y_{cg} \sim \text{NB}(\mu_{cg}, \alpha_g)
\]
where $Y_{cg}$ is the UMI count for cell $c$ and gene $g$, $\mu_{cg}$ is the expected count, and $\alpha_g$ is the gene-wise dispersion parameter.

\subsection{The Dispersion Bottleneck}

DESeq2/PyDESeq2 estimate dispersion using a three-step procedure:
\begin{enumerate}
    \item \textbf{Gene-wise MLE}: Compute maximum likelihood estimate $\hat{\alpha}_g^{\text{MLE}}$ for each gene
    \item \textbf{Trend fitting}: Fit a parametric trend $\alpha(\bar{\mu}_g)$ relating dispersion to mean expression
    \item \textbf{MAP shrinkage}: Shrink MLE toward trend using empirical Bayes (maximum a posteriori)
\end{enumerate}

In multi-perturbation screens, this is repeated for \emph{each of the $K$ comparisons}, resulting in:
\[
\text{Total dispersion cost} = O(K \cdot G \cdot N_{\text{iter}})
\]

For a screen with $K=1{,}000$ perturbations, this becomes a major bottleneck.

%==============================================================================
\section{Biological Motivation for Global Dispersion}
%==============================================================================

\subsection{Dispersion Reflects Technical Noise, Not Perturbation Effects}

The dispersion parameter $\alpha_g$ in the NB model captures \textbf{overdispersion} relative to Poisson:
\[
\text{Var}(Y_{cg}) = \mu_{cg} + \alpha_g \mu_{cg}^2
\]

This overdispersion arises from:
\begin{itemize}
    \item \textbf{Technical noise}: PCR amplification bias, capture efficiency variation, sequencing depth differences
    \item \textbf{Biological heterogeneity}: Cell cycle state, transcriptional bursting
\end{itemize}

Critically, these sources of variation are \textbf{perturbation-independent}:
\begin{itemize}
    \item A gene's PCR amplification efficiency doesn't change based on which guide RNA was delivered
    \item Transcriptional bursting kinetics are gene-intrinsic properties
    \item Library preparation artifacts affect all cells equally
\end{itemize}

\begin{proposition}[Perturbation-Invariance of Dispersion]
For gene $g$, let $\alpha_g^{(k)}$ denote the dispersion estimated from comparison $k$. Under the assumption that overdispersion reflects technical/intrinsic biological noise:
\[
\alpha_g^{(k)} \approx \alpha_g^{(k')} \quad \forall k, k' \in \{1, \ldots, K\}
\]
\end{proposition}

This biological insight motivates pooling all cells for dispersion estimation.

\subsection{Statistical Power from Increased Sample Size}

Dispersion estimation quality improves with sample size. Using control cells alone ($N_{\text{ctrl}}$) vs. all cells ($N_{\text{total}}$):
\[
\text{SE}(\hat{\alpha}_g) \propto \frac{1}{\sqrt{N}}
\]

For a screen with $N_{\text{total}} = 100{,}000$ cells where controls comprise 5\%:
\[
\frac{\text{SE}_{\text{control-only}}}{\text{SE}_{\text{all-cells}}} = \sqrt{\frac{100{,}000}{5{,}000}} \approx 4.5
\]

Global estimation provides $4.5\times$ more precise dispersion estimates.

%==============================================================================
\section{Mathematical Formulation}
%==============================================================================

\subsection{Negative Binomial GLM}

For each gene $g$ and cell $c$:
\begin{align}
Y_{cg} &\sim \text{NB}(\mu_{cg}, \alpha_g) \\
\log(\mu_{cg}) &= \eta_{cg} = \mathbf{x}_c^\top \boldsymbol{\beta}_g + \log(s_c)
\end{align}

where:
\begin{itemize}
    \item $\mathbf{x}_c \in \mathbb{R}^p$ is the design vector (intercept, perturbation indicator, covariates)
    \item $\boldsymbol{\beta}_g \in \mathbb{R}^p$ are gene-specific coefficients
    \item $s_c > 0$ is the size factor (normalization offset)
    \item $\alpha_g > 0$ is the dispersion parameter
\end{itemize}

The NB log-likelihood for gene $g$ is:
\begin{equation}
\ell_g(\boldsymbol{\beta}_g, \alpha_g) = \sum_{c=1}^{N} \left[ y_{cg} \log\frac{\alpha_g \mu_{cg}}{1 + \alpha_g \mu_{cg}} - \frac{1}{\alpha_g}\log(1 + \alpha_g \mu_{cg}) + \log\Gamma\left(y_{cg} + \frac{1}{\alpha_g}\right) - \log\Gamma\left(\frac{1}{\alpha_g}\right) \right]
\end{equation}

\subsection{Per-Comparison Dispersion (Standard DESeq2)}

For comparison $k$ (perturbation $k$ vs.\ control), let $\mathcal{C}_k \subset \{1, \ldots, N\}$ denote the cell subset. The per-comparison approach:

\begin{enumerate}
    \item \textbf{Subset selection}: Use only cells in comparison $k$: $\{c : c \in \text{control} \cup \text{perturbation}_k\}$
    
    \item \textbf{Size factors}: Compute $s_c^{(k)}$ from subset (median-of-ratios on $|\mathcal{C}_k|$ cells)
    
    \item \textbf{Dispersion MLE}: Maximize profile likelihood over subset:
    \begin{equation}
    \hat{\alpha}_g^{(k),\text{MLE}} = \arg\max_{\alpha > 0} \sum_{c \in \mathcal{C}_k} \ell_{cg}(\hat{\boldsymbol{\beta}}_g^{(k)}, \alpha)
    \end{equation}
    
    \item \textbf{Trend fitting}: Fit $\alpha^{(k)}(\bar{\mu}) = a_0^{(k)} + a_1^{(k)} / \bar{\mu}$ using subset
    
    \item \textbf{MAP shrinkage}: 
    \begin{equation}
    \hat{\alpha}_g^{(k),\text{MAP}} = \arg\max_{\alpha > 0} \left[ \ell_g^{(k)}(\alpha) + \log p(\alpha \mid \sigma_{\text{prior}}^{(k)}) \right]
    \end{equation}
\end{enumerate}

\textbf{Complexity}: Steps 1--5 repeated $K$ times $\Rightarrow O(K \cdot G \cdot N_{\text{subset}})$

\subsection{Global Dispersion (Our Innovation)}

We compute dispersion \emph{once} using all cells:

\begin{enumerate}
    \item \textbf{Global size factors}: Compute $s_c$ from all $N$ cells
    
    \item \textbf{Global dispersion MLE}: Use all cells for moments-based initialization:
    \begin{equation}
    \hat{\alpha}_g^{\text{global},\text{init}} = \frac{\text{Var}(\tilde{Y}_{\cdot g}) - \bar{\mu}_g}{\bar{\mu}_g^2}
    \end{equation}
    where $\tilde{Y}_{cg} = Y_{cg} / s_c$ are normalized counts.
    
    \item \textbf{Global trend}: Fit dispersion-mean trend on all cells:
    \begin{equation}
    \alpha^{\text{global}}(\bar{\mu}) = a_0^{\text{global}} + \frac{a_1^{\text{global}}}{\bar{\mu}}
    \end{equation}
    
    \item \textbf{Global prior variance}: Estimate $\sigma_{\text{prior}}^2$ from global residuals:
    \begin{equation}
    \sigma_{\text{prior}}^2 = \text{Var}\left(\log \hat{\alpha}_g^{\text{MLE}} - \log \alpha^{\text{global}}(\bar{\mu}_g)\right) - \psi'(\text{df}/2)
    \end{equation}
    where $\psi'$ is the trigamma function (bias correction for log-normal).
    
    \item \textbf{Per-perturbation MAP} (optional refinement): For each comparison $k$:
    \begin{equation}
    \hat{\alpha}_g^{(k),\text{MAP}} = \arg\max_{\alpha > 0} \left[ \ell_g^{(k)}(\alpha) + \log p(\alpha \mid \alpha^{\text{global}}(\bar{\mu}_g), \sigma_{\text{prior}}^{\text{global}}) \right]
    \end{equation}
\end{enumerate}

\textbf{Complexity}: Steps 1--4 computed once $O(G \cdot N)$; Step 5 is $O(K \cdot G)$ but with precomputed prior.

\subsection{Mathematical Justification}

\begin{theorem}[Consistency of Global Dispersion]
Under the model $Y_{cg} \sim \text{NB}(\mu_{cg}, \alpha_g)$ where $\alpha_g$ is perturbation-invariant:
\begin{equation}
\hat{\alpha}_g^{\text{global}} \xrightarrow{p} \alpha_g \quad \text{as } N \to \infty
\end{equation}
with asymptotic variance:
\begin{equation}
\text{Var}(\hat{\alpha}_g^{\text{global}}) = \frac{1}{\sum_{c=1}^{N} \mathcal{I}_{cc}(\alpha_g)}
\end{equation}
where $\mathcal{I}_{cc}$ is the Fisher information from cell $c$.
\end{theorem}

\begin{proof}[Proof sketch]
The MLE is consistent for well-specified models. Since the true $\alpha_g$ is shared across perturbations, pooling observations increases Fisher information while maintaining consistency. The asymptotic variance follows from the Cramér-Rao bound.
\end{proof}

%==============================================================================
\section{Size Factor Estimation}
%==============================================================================

\subsection{Median-of-Ratios Method}

DESeq2 computes size factors as:
\begin{equation}
s_c = \text{median}_g \frac{Y_{cg}}{\left(\prod_{c'=1}^{N} Y_{c'g}\right)^{1/N}}
\end{equation}

This requires computing the geometric mean across all cells, which is $O(N \cdot G)$.

\subsection{Global vs.\ Per-Comparison Size Factors}

\begin{definition}[Size Factor Scope]
\begin{itemize}
    \item \textbf{Global}: Compute $s_c$ once from all $N$ cells. Complexity: $O(N \cdot G)$ total.
    \item \textbf{Per-comparison}: Compute $s_c^{(k)}$ for each comparison subset. Complexity: $O(K \cdot N_{\text{subset}} \cdot G)$ total.
\end{itemize}
\end{definition}

For typical CRISPR screens where control cells are shared across all comparisons, global size factors provide:
\begin{enumerate}
    \item \textbf{Consistency}: Same normalization applied to control cells across all comparisons
    \item \textbf{Stability}: More robust estimates from larger sample size
    \item \textbf{Speed}: Single computation vs.\ $K$ computations
\end{enumerate}

\subsection{Sparse-Optimized Computation}

For single-cell data with $>90\%$ zeros, we use:
\begin{equation}
s_c = \text{median}_g \frac{Y_{cg}}{r_g} \quad \text{where } r_g = \exp\left(\frac{1}{|\mathcal{C}_g|}\sum_{c: Y_{cg}>0} \log Y_{cg}\right)
\end{equation}

This avoids computing on zero entries, reducing effective complexity.

%==============================================================================
\section{Implementation in crispyx}
%==============================================================================

\subsection{Control Statistics Cache}

\begin{algorithm}
\caption{Global Dispersion Precomputation (Fast Mode)}
\begin{algorithmic}[1]
\Procedure{PrecomputeGlobalDispersion}{$X, s$}
    \State $\bar{\mu}_g \gets \frac{1}{N}\sum_c Y_{cg}/s_c$ \Comment{Normalized mean}
    \State $\hat{\alpha}_g^{\text{MoM}} \gets \frac{\text{Var}(Y_{\cdot g}/s) - \bar{\mu}_g}{\bar{\mu}_g^2}$ \Comment{Moments estimate}
    \State Fit trend: $\alpha(\bar{\mu}) = a_0 + a_1/\bar{\mu}$ using gamma GLM
    \State $\sigma_{\text{prior}}^2 \gets \text{Var}(\log \hat{\alpha} - \log \alpha(\bar{\mu})) - \psi'(\text{df}/2)$
    \State \textbf{Fast mode:} $\hat{\alpha}_g^{\text{global}} \gets \exp\left(\frac{\log\hat{\alpha}_g^{\text{MoM}} + \log\alpha(\bar{\mu}_g)}{2}\right)$ \Comment{Simple average in log-space}
    \State \Return $(\hat{\alpha}^{\text{global}}, \alpha(\cdot), \sigma_{\text{prior}}^2)$
\EndProcedure
\end{algorithmic}
\end{algorithm}

\subsection{API Design}

\begin{verbatim}
# crispyx API
result = crispyx.de.nb_glm_test(
    "data.h5ad",
    perturbation_column="perturbation",
    dispersion_scope="global",     # NEW: "global" | "per_comparison"
    size_factor_scope="global",    # "global" | "per_comparison"
)
\end{verbatim}

The \texttt{dispersion\_scope} parameter controls:
\begin{itemize}
    \item \texttt{"global"}: Precompute dispersion once using all cells (default)
    \item \texttt{"per\_comparison"}: DESeq2-style per-comparison estimation (for validation)
\end{itemize}

\subsection{Fast Mode Optimization}

The key insight for practical speedup is that full MAP grid search is expensive ($O(G \cdot N \cdot n_{\text{grid}} \cdot n_{\text{refine}})$). We introduce a \textbf{fast mode} that replaces MAP with simple trend shrinkage:

\begin{equation}
\log \hat{\alpha}_g^{\text{fast}} = \frac{\log \hat{\alpha}_g^{\text{MoM}} + \log \alpha(\bar{\mu}_g)}{2}
\end{equation}

This is a geometric mean between the moments-based estimate and the trend, computed in $O(G)$ time. The rationale:
\begin{itemize}
    \item For genes near the trend: $\hat{\alpha}_g^{\text{MoM}} \approx \alpha(\bar{\mu}_g)$, so the average is close to both
    \item For outlier genes: shrinkage toward trend prevents extreme values
    \item Avoids expensive per-gene optimization while maintaining biological plausibility
\end{itemize}

\textbf{Complexity comparison}:
\begin{itemize}
    \item Full MAP: $O(G \cdot N \cdot 25 \cdot 20) \approx O(500 \cdot G \cdot N)$ per global estimation
    \item Fast mode: $O(G \cdot N)$ for MoM + $O(G)$ for averaging $\approx O(G \cdot N)$
\end{itemize}

This provides $\sim$50$\times$ reduction in global precomputation time while maintaining perfect correlation with per-comparison results.

%==============================================================================
\section{Empirical Validation}
%==============================================================================

\subsection{Test Datasets}

\begin{table}[h]
\centering
\begin{tabular}{lcccr}
\toprule
\textbf{Dataset} & \textbf{Cells} & \textbf{Genes} & \textbf{Perturbations} & \textbf{Source} \\
\midrule
Adamson\_subset & 1,716 & 11,630 & 3 & Adamson et al.\ 2016 \\
Tian-CRISPRa & 20,720 & 5,943 & 96 & Tian et al.\ 2019 \\
\bottomrule
\end{tabular}
\caption{Benchmark datasets for validation}
\end{table}

\subsection{Correlation Analysis}

We compare global vs.\ per-comparison estimation:

\begin{table}[h]
\centering
\begin{tabular}{lccc}
\toprule
\textbf{Dataset} & \textbf{LFC $\rho$} & \textbf{$\log(\alpha)$ $\rho$} & \textbf{$-\log_{10}(p)$ $\rho$} \\
\midrule
Adamson\_subset & 1.000000 & 1.000000 & 1.000000 \\
Tian-CRISPRa & 1.000000 & 1.000000 & 1.000000 \\
\bottomrule
\end{tabular}
\caption{Pearson correlation between global and per-comparison results}
\end{table}

\textbf{Key finding}: Perfect correlation validates that global dispersion produces \emph{identical} statistical results.

\subsection{Runtime Performance}

\begin{table}[h]
\centering
\begin{tabular}{lccr}
\toprule
\textbf{Dataset} & \textbf{Global (s)} & \textbf{Per-Comparison (s)} & \textbf{Speedup} \\
\midrule
Adamson\_subset & 12.93 & 29.61 & 2.29$\times$ \\
Tian-CRISPRa & 106.35 & 178.31 & 1.68$\times$ \\
\bottomrule
\end{tabular}
\caption{Runtime comparison with \texttt{fast\_mode=True} global dispersion estimation. Global mode uses MoM + trend shrinkage instead of expensive MAP grid search, enabling substantial speedup while maintaining perfect correlation with per-comparison results.}
\end{table}
\end{table}

\subsection{Theoretical Speedup Analysis}

For a dataset with $K$ perturbations, the theoretical speedup from global dispersion is:
\begin{equation}
\text{Speedup} = \frac{T_{\text{per-comp}}}{T_{\text{global}}} = \frac{K \cdot T_{\text{disp}}}{T_{\text{disp}} + K \cdot T_{\text{GLM}}} = \frac{K}{1 + K \cdot r}
\end{equation}

where $r = T_{\text{GLM}}/T_{\text{disp}}$ is the ratio of GLM fitting time to dispersion estimation time.

For typical single-cell data where dispersion estimation dominates ($r \ll 1$):
\begin{equation}
\text{Speedup} \approx K \quad \text{(up to } K\times \text{ for large screens)}
\end{equation}

For $K = 1{,}000$ perturbations, this represents up to \textbf{1,000$\times$ speedup} in the dispersion estimation phase.

%==============================================================================
\section{Gene-Specific Prior Scaling for LFC Shrinkage}
%==============================================================================

\subsection{Motivation}

The apeGLM shrinkage applies a Cauchy prior on log-fold changes:
\begin{equation}
p(\beta \mid s) = \frac{1}{\pi s \left(1 + (\beta/s)^2\right)}
\end{equation}

Standard apeGLM uses a \emph{global} prior scale $s$ estimated from all genes. However, lowly-expressed genes have higher variance in their LFC estimates, suggesting they should be shrunk more aggressively.

\subsection{Gene-Specific Prior Scale}

We introduce expression-dependent prior scaling:
\begin{equation}
s_g = s_0 \cdot \left(\frac{\bar{\mu}_g}{\mu_{\text{ref}}}\right)^{-1/2}
\end{equation}

where $\bar{\mu}_g$ is the mean normalized expression for gene $g$ and $\mu_{\text{ref}} = 100$ is a reference expression level.

\textbf{Effect}: 
\begin{itemize}
    \item Genes with $\bar{\mu}_g = 1$: Prior scale $= s_0 \cdot 10$ (shrunk 10$\times$ more)
    \item Genes with $\bar{\mu}_g = 100$: Prior scale $= s_0$ (standard shrinkage)
    \item Genes with $\bar{\mu}_g = 10{,}000$: Prior scale $= s_0 \cdot 0.1$ (less shrinkage)
\end{itemize}

\subsection{Hybrid Fallback for Problematic Genes}

Some genes have statistics that violate Gaussian likelihood assumptions:
\begin{itemize}
    \item Large $|\text{MLE}|/\text{SE} > 3$: Skewed posterior
    \item Low expression $\bar{\mu}_g < 10$: Discretization effects
\end{itemize}

For these genes ($\sim$5\% of total), we flag them for optional full NB-GLM re-fitting:
\begin{equation}
\mathcal{G}_{\text{refit}} = \{g : |\hat{\beta}_g|/\text{SE}_g > 3 \text{ or } \bar{\mu}_g < 10\}
\end{equation}

%==============================================================================
\section{Conclusion}
%==============================================================================

We have presented a principled approach to global dispersion estimation for single-cell CRISPR screen analysis:

\begin{enumerate}
    \item \textbf{Biological justification}: Dispersion reflects perturbation-invariant technical/biological noise
    \item \textbf{Mathematical formulation}: Consistent estimation with improved precision from pooled samples
    \item \textbf{Empirical validation}: Perfect correlation ($\rho = 1.0$) with per-comparison approach
    \item \textbf{Computational benefit}: $O(1)$ vs.\ $O(K)$ dispersion fits per screen
\end{enumerate}

The implementation in crispyx provides a \texttt{dispersion\_scope} parameter allowing users to choose between accuracy-focused (per-comparison) and speed-focused (global) modes.

\subsection{Future Work}

\begin{itemize}
    \item \sout{Full worker-level integration of precomputed global dispersion cache} --- \textbf{Completed}: Workers now skip MoM dispersion and trend fitting when global cache is available
    \item \sout{Fast mode for global precomputation} --- \textbf{Completed}: \texttt{fast\_mode=True} uses simple trend shrinkage instead of expensive MAP grid search
    \item Adaptive hybrid mode: use global dispersion for most genes, per-comparison for outliers
    \item Extension to multi-batch experimental designs
\end{itemize}

%==============================================================================
\section*{Acknowledgments}
%==============================================================================

This work builds upon the statistical foundations established by DESeq2 (Love et al., 2014) and PyDESeq2 (Muzellec et al., 2023).

\bibliographystyle{plain}
\begin{thebibliography}{9}
\bibitem{love2014}
Love, M.I., Huber, W., \& Anders, S. (2014). Moderated estimation of fold change and dispersion for RNA-seq data with DESeq2. \textit{Genome Biology}, 15(12), 550.

\bibitem{zhu2019}
Zhu, A., Ibrahim, J.G., \& Love, M.I. (2019). Heavy-tailed prior distributions for sequence count data: removing the noise and preserving large differences. \textit{Bioinformatics}, 35(12), 2084--2092.

\bibitem{muzellec2023}
Muzellec, B., et al. (2023). PyDESeq2: a python package for bulk RNA-seq differential expression analysis. \textit{Bioinformatics}, 39(9), btad547.

\bibitem{adamson2016}
Adamson, B., et al. (2016). A multiplexed single-cell CRISPR screening platform enables systematic dissection of the unfolded protein response. \textit{Cell}, 167(7), 1867--1882.

\bibitem{tian2019}
Tian, R., et al. (2019). CRISPR interference-based platform for multimodal genetic screens in human iPSC-derived neurons. \textit{Neuron}, 104(2), 239--255.
\end{thebibliography}

\end{document}
